% !PRINTED 5
\weberpage{5}

Bei den im ersten Bande betrachteten Permutationsgruppen wird man leicht die Merkmale des allgemeinen Gruppenbegriffes erkennen.

Wir ziehen nun aus dieser Definition zunächst einige ganz allgemeine Folgerungen.

Nach 4. giebt es für jedes gegebene $b$ \gesperrt{ein} Element $e$ in $P$, das der Bedingung
\weq{1}{eb = b}
genügt, und dies $e$ ist von $b$ unabhängig; denn aus (1) folgt für jedes $c$
\[
ebc = bc,
\]
und $bc$ kann nach 4. jedes Element in $P$ bedeuten. Ebenso giebt es ein Element $e'$, das für jedes $b$ der Bedingung
\weq{2}{be' = b}
genügt. Dies Element $e'$ ist aber von $e$ nicht verschieden; denn setzen wir $b = e'$ in (1) und $b = e$ in (2), so folgt
\[
e e' = e', \, . \, e e' = e, % [sic: the source prints a full stop between the two identities]
\]
also
\[
e = e'.
\]

Das Element $e$ ändert nichts, wenn es mit irgend welchen Elementen aus $P$ componirt wird, und wird die \gesperrt{Einheit der Gruppe} genannt.

In vielen Fällen kann es ohne Missverständniss geradezu mit „1“ bezeichnet werden.

Zu jedem Element $a$ giebt es nach 4. ein bestimmtes Element $a^{-1}$, das der Bedingung
\weq{3}{a^{-1} a = e}
genügt. Aus (3), (1) und (2) folgt
\[
a^{-1} a a^{-1} = e a^{-1} = a^{-1} = a^{-1} e,
\]
und folglich nach 3.
\weq{4}{a a^{-1} = e.}

Die beiden Elemente $a$, $a^{-1}$ heissen zu einander \gesperrt{entgegengesetzt} oder \gesperrt{reciprok}.

In besonderen Fällen kann bei der Composition der Elemente einer Gruppe $P$ auch das commutative Gesetz gelten, d.\,h. es kann für je zwei Elemente $a$, $b$ der Gruppe
\[
a b = b a
\]
sein.
